% !TEX program = lualatex
% LTeX: language=it

%%% 
%%% Usare il compositore «LuaLaTeX»
%%% 

\documentclass[11pt]{beamer}


\makeatletter%
\newcommand\controllaClasse[2]{%
    \@ifclassloaded{article}{#1}{#2}%
}%
\makeatother%

% \controllaClasse{% Classe «article»
% }{% Classe «beamer»
% }%

\controllaClasse{% Classe «article»
    \usepackage[%
        % showframe,%
        % showcrop,%
        hmargin=1cm,%
        bottom=2cm,%
        top=2cm,%
        twocolumn%
    ]{geometry}%
    % \usepackage{multicol} % Pacchetto per gestire le colonne, specie quando si ha un documento misto
    \setlength{\columnsep}{0.75cm}%
}{% Classe «beamer»
    % \usepackage{showframe}
    \setlength\parskip{7.5pt} % Impone la distanza tra due capoversi uguale a «7.5pt» ove il valore di base nella classe Beamer è «0pt»
    % \setlength{\parindent}{15pt} % Impone il rientro al valore di base 15pt, poiché la classe Beamer la pone nulla se non altrimenti specificato
}%

\controllaClasse{% Classe «article»
    \usepackage[british]{babel}
    % \usepackage[utf8]{inputenc} % Se si sta lavorando con LuaLaTeX o XeLaTeX allora si può commentare mentre con pdfLaTeX è necessario    
}{% Classe «beamer»
    \usepackage[italian]{babel}
}%

\usepackage[T1]{fontenc}
% \usepackage[T1]{tipa}
\usepackage{microtype} % Migliora la giustificazione del testo

% Cambia il fonte usato nel documento secondo le filze .ttf salvate in «path»
\usepackage{fontspec} % È necessario il compilatore «LuaLaTeX» o «XeLaTeX»
\controllaClasse{%
    \setmainfont{Vollkorn}[%
        Path = ../Riferimenti/,	        % Percorso
        Extension = .ttf,               % Estensione
        UprightFont = *-Regular,        % Fonte tondo
        ItalicFont = *-Italic,          % Fonte corsivo
        BoldFont = *-Bold,              % Fonte grassetto
        BoldItalicFont = *-BoldItalic   % Fonte grassetto corsivato
    ]%
}{%
    \setsansfont{Vollkorn}[%
        Path = ../Riferimenti/,	        % Percorso
        Extension = .ttf,               % Estensione
        UprightFont = *-Regular,        % Fonte tondo
        ItalicFont = *-Italic,          % Fonte corsivo
        BoldFont = *-Bold,              % Fonte grassetto
        BoldItalicFont = *-BoldItalic   % Fonte grassetto corsivato
    ]%
}%
% Le quattro filze per il fonte sono salvati in 
% «./Percorso/della/Filza/Fonte-Regular.ttf»
% «./Percorso/della/Filza/Fonte-Italic.ttf»
% «./Percorso/della/Filza/Fonte-Bold.ttf»
% «./Percorso/della/Filza/Fonte-BoldItalic.ttf»

\usepackage[backend=biber,style=numeric-comp]{biblatex}%
\addbibresource{../Riferimenti/Bibliografia.bib}
\nocite{*} % Cite all entries inside the «.bib» file
\usepackage{csquotes}

% Pacchetti per vari comandi matematici
\usepackage{amsmath,
            amsfonts,
            amsthm,
            amssymb,
            amsthm}

\usepackage{bbold}     % Per es. l'1 con grassetto da lavagna
\usepackage{mathtools} % Per es. \xrightleftharpoons
\usepackage{empheq}
\usepackage{cancel}

\usepackage{nicematrix,
            array}

\usepackage{graphicx} % Pacchetto per le immagini
\usepackage[shortlabels]{enumitem} % Pacchetto per gestire le liste
\setlist[itemize]{label=$\diamond$}

\usepackage{tikz}
\usetikzlibrary{
    tikzmark,
    positioning,
    arrows.meta,
    fit,
    calc,
    decorations.pathreplacing,
    calligraphy,
    spy,
}
\usepackage{pgfplots}
\usepgfplotslibrary{groupplots}
\pgfplotsset{compat=1.18}

\usepackage{caption,    % Pacchetto per le didascalie dentro le minipagine
            subcaption} % Pacchetto per avere due figure affiancate
\usepackage{float,
            stfloats}

\usepackage{xcolor}   % Pacchetto per colorare il testo
\usepackage{hyperref} % Pacchetto per i collegamenti ipertestuali
\newcommand{\teqref}[2]{\hyperref[#1]{#2}} % Riferimenti applicati a un testo

% Il comando da usare per questo pacchetto è \cref{} anziché \eqref{} o \ref{}
\controllaClasse{% Classe «article»    
    \usepackage{cleveref} % Pacchetto per citare multiple equazioni
    %
    \crefrangeformat{equation}{(#3#1#4--#5#2#6)}%
    \crefmultiformat{equation}{(#2#1#3}{, #2#1#3)}{, #2#1#3}{, #2#1#3)}%
    \crefrangemultiformat{equation}{(#3#1#4--#5#2#6}{, #3#1#4--#5#2#6)}{, #3#1#4--#5#2#6}{, #3#1#4--#5#2#6)}%
    %
    \crefrangeformat{enumi}{#3#1#4--#5#2#6}%
    \crefmultiformat{enumi}{#2#1#3}{ and #2#1#3}{, #2#1#3}{ and #2#1#3}%
    %
    \crefrangeformat{figure}{#3#1#4--#5#2#6}%
    \crefmultiformat{figure}{#2#1#3}{ and #2#1#3}{, #2#1#3}{ and #2#1#3}%
    \crefrangemultiformat{figure}{#3#1#4--#5#2#6}{ and #3#1#4--#5#2#6}{, #3#1#4--#5#2#6}{ and #3#1#4--#5#2#6}%
    %
    \crefrangeformat{section}{#3#1#4--#5#2#6}%
    \crefmultiformat{section}{#2#1#3}{ and #2#1#3}{, #2#1#3}{ and #2#1#3}%
}{% Classe «beamer»
    \usepackage{cleveref} % Pacchetto per citare multiple equazioni
    %
    \crefrangeformat{enumi}{#3#1#4--#5#2#6}%
    \crefmultiformat{enumi}{#2#1#3}{ e #2#1#3}{, #2#1#3}{ e #2#1#3}%
}%

\usepackage{siunitx}
\sisetup{
    inter-unit-product=\ensuremath{{}\cdot{}},
    scientific-notation=false,
}

\controllaClasse{}{% Classe «beamer»
    \beamertemplatenavigationsymbolsempty % Elimina la barra di navigazione in fondo alle diapo
    %
    \usepackage{ragged2e} % Pacchetto per giustificare il testo
    \justifying%
    %
    \renewcommand{\implies}{\DOTSB\Longrightarrow} % Nella definizione standarda «\implies» ha forma «\DOTSB \;\Longrightarrow \;», ove «\DOTSB» corrisponde a «\relax» vale a dire che non fa niente; tuttavia beamer [mannaggia non so perché] frigna che i due «\;» sono invalidi (in generale in modalità matematica par essere cosí dato che non li accetta), quindi è necessario ridefinire «\implies» senza i due simboli cosí da non avere errori imprevisti
    %
    % Settaggio delle didascalie
    % \captionsetup{aboveskip=5pt plus 0pt minus 0pt,           % Spazio sopra la didascalia
    %               belowskip=-2.5pt plus 2.5pt minus 2.5pt}    % Spazio sotto la didascalia
}%
%     %%% Ridefinizione di \@maketitle
\makeatletter % Serve solo per non far frignare LaTeX per il simbolo «@»
\def\@maketitle{
\newpage

% la macro «\null» aggiunge solo una scatola vuota orizzontale ed è responsabile dello spazio esagerato bianco prima di \maketitle, per cui commentarlo elimina il problema estetico.
% Cosí agendo lo spazio verticale prima dello spazio è interamente determinato dall'impaginazione del documento, perciò se si vuole aggiungere spazio bisogna farlo col pacchetto geometry
% Ho anche notato che senza «\null» il comando «\vskip» non ha alcun effetto
% \null
% \vskip 2em

\begin{center}
    \let\footnote\thanks
    {\LARGE\@title\par}
    \vskip 1em
    {\large 
        \lineskip .5em
        \begin{tabular}[t]{c}
            \@author 
        \end{tabular}\par}%
    % \vskip 0em
    % {\large \@date }
\end{center}%
% \par
\vskip -1em
}
\makeatother % Serve solo per non far frignare LaTeX per il simbolo «@»
\newcommand{\derT}[3][d]{\frac{\mathrm{#1}#2}{\mathrm{#1}#3}} % Derivata totale
\newcommand{\derTiL}[3][d]{\mathrm{#1}#2/\mathrm{#1}#3}       % Derivata totale in linea
\newcommand{\derP}[2]{\frac{\partial#1}{\partial#2}}	      % Derivata parziale

\DeclareMathOperator\VolOp{Vol}
\newcommand{\Vol}[1]{\VolOp\left(#1\right)}

\newcommand\de{\mathrm{d}}
\DeclareMathOperator\diag{diag}
\DeclareMathOperator\trace{tr}

\newcommand{\SpazMate}[3]{% Spaz[iatura] Mate[matica]
    \thinmuskip=#1mu\relax
    \medmuskip=#2mu\relax
    \thickmuskip=#3mu\relax
}
\SpazMate{1}{2}{3} % Seleceting math spacing

\newcommand\DeltaD{\Delta^{\mathbf D}}

\newcommand\RPlus{\mathbb R^+}
\newcommand\RPlusStar{\mathbb R_*^+}
\newcommand\NPlus{\mathbb N^+}


\newcommand\FKdiscrete{\teqref{eqFKDiscrete}{[FK]$_d$}}
\newcommand\HDdiscrete{\teqref{eqHDDiscrete}{[HD]$_d$}}
\newcommand\SMOdiscrete{\teqref{eqSMODiscrete}{[SMO]$_d$}}

\newcommand\FKcontinuous{\teqref{eqFKContinuous}{[FK]$_c$}}
\newcommand\HDcontinuous{\teqref{eqHDContinuous}{[HD]$_c$}}
\newcommand\SMOcontinuous{\teqref{eqSMOContinuous}{[SMO]$_c$}}

\newcommand{\TotM}{M_{\text{tot}}}
\newcommand{\TotP}{P_{\text{tot}}}
\newcommand\bOne{\boldsymbol1}

\newcommand\mInf{m_\infty}
\newcommand\MInf{M_\infty}
\newcommand\PInf{P_\infty}

\newcommand\Mass{\text{M}}
% \newcommand\Year{\text{a}}
\DeclareSIUnit\year{a}
% \newcommand\Molarity{\text{M}}
\DeclareSIUnit\molarity{M}

\newcommand\selezionaLingua[2]{%
  \iflanguage{italian}{#1}% Se la lingua corrente è l'inglese usa #1
  {% Altrimenti se la lingua è in italiano usa #2
    \iflanguage{british}{#2}%
    {%
      Linguaggio non considerato % Se non è né italiano né inglese restituisce inglese
    }%
  }%
}

\newcommand\didascalia[1]{%
    \captionsetup{aboveskip=-.2cm}
    \addcontentsline{lof}{figure}{\protect\numberline{\thefigure}#1}%
    \caption*{\figurename~\thefigure: #1}
}

\newcommand\sottoDidascalia[1]{%
    \refstepcounter{subfigure}%
    (\thesubfigure) #1%
}

\newcommand\dimTesto[1]{%
    \pgfmathparse{#1*1.2}%
    \fontsize{#1 pt}{\pgfmathresult pt}\selectfont%
}

\makeatletter
\newcommand{\includetikzgroupplotfigure}
    {\@ifstar{\includetikzgroupplotfigure@star}
    {\includetikzgroupplotfigure@nostar}}

\newcommand{\includetikzgroupplotfigure@nostar}[5]{%
    \begin{figure}[#1]%
        \centering%
        \refstepcounter{figure}%
        \resizebox{#2}{!}{#3}%
        \didascalia{#4}%
        \label{#5}%
    \end{figure}%
}

\newcommand{\includetikzgroupplotfigure@star}[6]{%
    \begin{figure*}[#1]%
        \vspace{#2}%
        \centering%
        \refstepcounter{figure}%
        \resizebox{#3}{!}{#4}%
        \didascalia{#5}%
        \label{#6}%
    \end{figure*}%
}
\makeatother

\newcommand\ie{\textit{i.e.}}

\newcommand\Brain{\mathcal B}
\newcommand\STCyl{\mathcal Q} % Space-time cylinder

\newcommand\vColla[2]{\vspace{#1 plus #2 minus #2}}
\newcommand\hColla[2]{\hspace{#1 plus #2 minus #2}}

\newcommand\HadProd{\smash{\vcenter{\hbox{\scalebox{0.7}{$\odot$}}}}}

\newcommand\HeII{\mathrm{He}_2}

\begin{document}

    \title{{%
        \makebox[.9\textwidth]{% Non so perché se scrivo «\makebox[.9\textwidth]» il testo è centrato rispetto alla diapo mentre con «\makebox[\textwidth]» non lo è
            \parbox{1.125\textwidth}{\centering%
                \LARGE\emph{\textbf{%
                    Analisi di modelli neurali su rete per la malattia di Alzheimer tramite la propagazione di proteine tau malpiegate similprioniche %https://www.treccani.it/vocabolario/simil/
                }}%
        }}}\texorpdfstring{\\\vspace{10pt}}{}%
        {\large Corso di Modelli Matematici per la Biomedicina}
    }%
    \author{Taralli Valerio (s332689)\texorpdfstring{\\}{} Cappadonia Silvia (s329228)}
    \date{}%04/09/2024}

    \maketitle
    
    %%%
    %%% Premessa
    %%%
    	
    \begin{frame}{Premessa}
        Questa è la presentazione ...

        Si divide in tre parti
    \end{frame}

    
    %%%
    %%% 1° Parte - Introduzione
    %%%
    
    \part{Introduzione}
    \begin{frame}\partpage\end{frame}


    \section{Fenomenologia}
    
        \subsection{Ipotesi prionica}
        
            \begin{frame}\frametitle{\insertsection{} - \insertsubsection{}}
                H1 H2 H3 ...
            \end{frame}
    
    %%%
    %%% 2° Parte - Modelli cinetici
    %%%

    \part{Modelli cinetici}
    \begin{frame}\partpage\end{frame}


    \section{Premessa}
        
        \subsection{Notazione}

            \begin{frame}\frametitle{\insertsection{} - \insertsubsection{}}
                Cilindro spazio-temporale

                Rpiú ed RpiúStella
            \end{frame}


    \section{Modelli continui}
    
        \subsection{Eterodimero (ED)}

            \begin{frame}\frametitle{\insertsection{} - \insertsubsection{}}
                [HD]$_c$
            \end{frame}

        \subsection{Fisher-Kolmogorov (FK)}
        
            \begin{frame}\frametitle{\insertsection{} - \insertsubsection{}}
                [FK]$_c$
            \end{frame}

        \subsection{Smoluchowski (SMO)}

            \begin{frame}\frametitle{\insertsection{} - \insertsubsection{}}
                Forma originale
            \end{frame}

            \begin{frame}\frametitle{\insertsection{} - \insertsubsection{}}
                Ipotesi semplificative
            \end{frame}

            \begin{frame}\frametitle{\insertsection{} - \insertsubsection{}}
                [SMO]$_c$
            \end{frame}


    \section{Modelli discreti}

        \subsection{Discretizzazione spaziale}

            \begin{frame}\frametitle{\insertsection{} - \insertsubsection{}}
                Logica e forma con rho
            \end{frame}

            \begin{frame}\frametitle{\insertsection{} - \insertsubsection{}}
                Modelli discreti: FK, ED e SMO
            \end{frame}


        \subsection{Tempi caratteristici}

            \begin{frame}\frametitle{\insertsection{} - \insertsubsection{}}
                Logica
            \end{frame}

            \begin{frame}\frametitle{\insertsection{} - \insertsubsection{}}
                Definizione di $\tau_d$ e di $\tau_g$ (tralascia le dimostrazioni)
            \end{frame}

            \begin{frame}\frametitle{\insertsection{} - \insertsubsection{}}
                Classificazione delle dinamiche
            \end{frame}

            \begin{frame}\frametitle{\insertsection{} - \insertsubsection{}}
                Immagine dei confronti (Forse da unire colla diapo precedente?)
            \end{frame}

       
        \subsection{Discretizzazione temporale}

            \begin{frame}\frametitle{\insertsection{} - \insertsubsection{}}
                Logica e forma generale esplicita e implicita
            \end{frame}


    \section{Analisi asintotica}

        \subsection{Modello FK}

            \begin{frame}\frametitle{\insertsection{} - \insertsubsection{}}
                Logica molto veloce
            \end{frame}

        \subsection{Modello ED}

            \begin{frame}\frametitle{\insertsection{} - \insertsubsection{}}
                EP con analisi (senza dimostrazione, specie perché nell'appendice)
            \end{frame}

        \subsection{Modello SMO}

            \begin{frame}\frametitle{\insertsection{} - \insertsubsection{}}
                Tre equazioni principali (senza dimostrazione, però magari accenna alla somma)
            \end{frame}

        \subsection{Tempo di convergenza}

            \begin{frame}\frametitle{\insertsection{} - \insertsubsection{}}
                Formula empirica
            \end{frame}

            \begin{frame}\frametitle{\insertsection{} - \insertsubsection{}}
                Osservazioni piú immagine (potrebbe occupare piú diapo)
            \end{frame}


    %%%
    %%% 3° Parte - Risultati e confronti
    %%%


    \part{\texorpdfstring{Risultati\\\&\\Confronti}{Risultati \& Confronti}}
    \begin{frame}\partpage\end{frame}

    
    \section{Complicazioni}

        \subsection{\texorpdfstring{$\mathbf L$}{L} dinamica}

            \begin{frame}\frametitle{\insertsection{} - \insertsubsection{}}
                Logica
            \end{frame}
        
            \begin{frame}\frametitle{\insertsection{} - \insertsubsection{}}
                Classificazione lati
            \end{frame}

            \begin{frame}\frametitle{\insertsection{} - \insertsubsection{}}
                Funzioni di decadimento
            \end{frame}

            \begin{frame}\frametitle{\insertsection{} - \insertsubsection{}}
                Ridefinizione della $\mathbf L$
            \end{frame}

        \subsection{Accumulazione}

            \begin{frame}\frametitle{\insertsection{} - \insertsubsection{}}
                Logica
            \end{frame}

        \subsection{Dipendenza dalla dimensione}

            \begin{frame}\frametitle{\insertsection{} - \insertsubsection{}}
                Logica della diffusione e della depurazione
            \end{frame}
            
    
    \section{Impostazioni}

        \subsection{Premessa}

            \begin{frame}\frametitle{\insertsection{} - \insertsubsection{}}
                Scelta della rho generale e dei nodi infetti
            \end{frame}    
        
        \subsection{Modello FK}

            \begin{frame}\frametitle{\insertsection{} - \insertsubsection{}}
                Parametri e CI
            \end{frame}    
        
        \subsection{Modello ED}

            \begin{frame}\frametitle{\insertsection{} - \insertsubsection{}}
                Parametri e CI
            \end{frame}    
        
        \subsection{Modello ED}

            \begin{frame}\frametitle{\insertsection{} - \insertsubsection{}}
                Parametri e CI
            \end{frame}    
        

        \subsection{Avvertenza}

            \begin{frame}\frametitle{\insertsection{} - \insertsubsection{}}
                Avvertenza sulla natura teorica dei parametri e sulla difficile riproduzione
            \end{frame}
    

    \section{Risultati}

        \subsection{Modelli FK ed ED}

            \begin{frame}\frametitle{\insertsection{} - \insertsubsection{}}
                Immagini
            \end{frame}    
        
            \begin{frame}\frametitle{\insertsection{} - \insertsubsection{}}
                Lunghe spiegazioni
            \end{frame}            

        \subsection{Modello SMO}

            \begin{frame}\frametitle{\insertsection{} - \insertsubsection{}}
                Immagini
            \end{frame}    
        
            \begin{frame}\frametitle{\insertsection{} - \insertsubsection{}}
                Lunghe spiegazioni
            \end{frame}    


    %%%
    %%% 4° Parte - Conclusioni
    %%%


    \part{Conclusioni}
    \begin{frame}\partpage\end{frame}

        \begin{frame}\frametitle{\insertpart{}}
            Commenti sui risultati
        \end{frame}    

    \setbeamertemplate{part page}{
        \vbox to\paperheight{%
            \vfil
            \begingroup
                \centering
                \begin{beamercolorbox}[
                    sep=16pt,
                    center,
                    wd=\textwidth,
                    colsep=-4bp,
                    rounded=true,
                    shadow=\beamer@themerounded@shadow
                ]{part title}
                \usebeamerfont{part title}\insertpart\par
                \end{beamercolorbox}
            \endgroup
            \vfil
        }
    }

    \part{Grazie dell'attenzione}
    \begin{frame}\partpage\end{frame}
        
    % \part{Riferimenti}

    % \begin{frame}[allowframebreaks,t]\frametitle{\insertpart{}}
    %     \printbibliography
    %     % \setbeamertemplate{bibliography item}{\insertbiblabel}
    %     % \bibliographystyle{ieeetr}
    %     % \bibliography{Bibliografia.bib}
    % \end{frame}

\end{document}

    % \begin{frame}\frametitle{\insertsection{} - \insertsubsection{}}
    % \end{frame}    
